\section{Related work}
\label{sec:related}
\paragraph{Fewer items inside one benchmark.}
tinyBenchmarks, anchor points, metabench and later work shrink the item set of a benchmark and estimate that benchmark's score from the retained items \citep{polo2024tinybenchmarks,vivek2024anchor,kipnis2024metabench,perlitz2024efficient,hofmann2025fluid,yauney2025reliable}.
Their unit of cost is an item; ours is a whole benchmark, and the quantity to infer is the score on the \emph{other} benchmarks.
\paragraph{Agreement and low-dimensional predictors.}
BenchBench quantifies how benchmarks agree \citep{perlitz2024benchbench}; observational scaling laws, Sloth and a latent-variable scaling law predict benchmark scores through a few capability factors estimated from many public models, from model features rather than from benchmarks already run \citep{ruan2024observational,polo2024sloth,cai2025latent}.
Factor analyses of language-model scores find a few broad abilities \citep{burnell2023revealing,ilic2023evidence}; Burnell et al.\ suggest that benchmarks could be streamlined to one task per ability.
We take that low-dimensional structure as our premise and turn it into a selection of benchmarks, a held-out prediction and a derived number $k$ to run.
\paragraph{Item response theory and adaptive testing.}
Psychometric adaptive testing selects items by their information about a latent ability \citep{segall1996multidimensional}, and item response models have been used to build NLP evaluation scales, compare test sets and reweight evaluation examples \citep{lalor2016building,vania2021comparing,rodriguez2021evaluation,lalor2024irt}; recent work argues for the same approach for language models and applies Fisher-information item selection inside one benchmark \citep{zhuang2023position,li2025atlas}.
Their unit is again the item and the target is the ability or the score of the benchmark it comes from; we choose among whole benchmarks, with their unequal cost, to predict the scores on those not run.
\paragraph{Cold-start elicitation and experimental design.}
Choosing which items a new user should rate has the structure of our problem \citep{golbandi2011adaptive,elahi2016survey}; the machinery for choosing observations in a Gaussian model is greedy sensor placement with submodularity guarantees \citep{krause2005nonmyopic,krause2008nearoptimal}, and for variance reduction the submodularity ratio of \citet{das2011submodular}, see also \citet{khanna2017scalable}.
Column subset selection and leverage scores choose representative columns \citep{cohen2017input}; matrix completion recovers entries observed at random \citep{candes2009exact,mazumder2010spectral}, whereas here the observed set is chosen, budgeted, and lies in a new row.
We use this machinery and claim no novelty for it.
